% Scope of this working report; administrative fields are intentionally unfilled.
\FrontChapter{RESEARCH PROGRESS STATUS}
\noindent\textbf{Evidence cutoff: 7 November 2026.} This is a research progress report within the
existing thesis document. Its purpose is to explain the problem, the decisions made, the experiments
completed and the questions still open. The official project title and Faculty template are retained;
unfilled defence, committee and acceptance fields do not represent completed academic procedures.

\begin{description}
  \item[Verified observations.] Public reporter-assay labels, recorded configurations, saved
  predictions and reproducible calculations support the reported development results. Completion
  of the local runnable graph means 633 jobs completed, with two extensions awaiting reviewed inputs.
  It does not establish independent biological validation or acceptance of a final model.
  \item[Evidence boundaries.] The restricted 91-peptide benchmark, the full-PSI 235-peptide development
  population and the locked same-laboratory transfer population answer different questions. Their
  metrics are reported separately. Three seeds sharing an outer partition are repeated fits rather
  than three independent test cohorts.
  \item[Interpretation and hypotheses.] An observed error reduction is distinguished from a proposed
  biological or statistical explanation. A same-budget retrospective OOF replay now supports
  improved positive-hit prioritisation over W5-D and random. It is not a new biological experiment:
  prospective yield, laboratory time/cost savings and clinical benefit remain unmeasured.
  \item[Open work.] Independent source validation, LoRA transfer, final serving refit/calibration/export
  and expanded software acceptance remain open. Missing information is classified explicitly in
  the accompanying review record instead of being replaced with assumed values.
\end{description}

The introduction states the current research questions. Chapter~\ref{chap:background} explains the
necessary biological and computational concepts. Section~\ref{sec:research-progress} reconstructs
the research decisions and data versions. Chapter~\ref{chap:evaluation} preserves unsuccessful
controls alongside the new results and assesses what each question can currently support.
The conclusion therefore records provisional findings and the next evidence needed.
